% Part: first-order-logic
% Chapter: introduction
% Section: formulas

\documentclass[../../../include/open-logic-section]{subfiles}

\begin{document}

\olfileid{fol}{int}{fml}

\section{\usetoken{P}{formula}}

మొదటిస్థాయి విధేయ తర్కపు \tetoken{వాక్యాలను}{!!{sentence}s} కఠినంగా నిర్దేశించడానికి,
\tetoken{చరాలు}{!!{variable}s} వాడకం వల్ల వచ్చే సమస్యలను పరిష్కరించడానికి మనం
అనుసరించే పద్ధతి ఇదీ. ముందుగా \emph{వేరొక} వ్యక్తీకరణల సమితిని,
అంటే \tetoken{సూత్రాలను}{!!{formula}s} నిర్వచిస్తాం. ఆ తరువాత వాటిలో \tetoken{చరాలు}{!!{variable}s}
పోషించే పాత్రను పరిశీలించవచ్చు; ``స్వేచ్ఛా,'' ``బద్ధ''
\tetoken{చరాలు}{!!{variable}s} అనే మరికొన్ని భావనల ఆధారంగా \tetoken{ఒక వాక్యం}{!!a{sentence}} అంటే ఏమిటో
నిర్వచించవచ్చు (అంటే, స్వేచ్ఛా \tetoken{చరాలు}{!!{variable}s} లేని \tetoken{ఒక సూత్రం}{!!a{formula}}).
``\tetoken{వాక్యం}{!!{sentence}}'' నిర్వచనాన్ని నిర్వహించడానికి ఇది సులభమవుతుందనే
కారణంతో మాత్రమే ఇలా చేయడం లేదు; సంతృప్తి అనే అర్థపర భావనను మనం
నిర్వచించే విధానానికి కూడా ఇది కీలకం అవుతుంది.

ఒకే \tetoken{విధేయ సంకేతం}{!!{predicate}}~$\Obj P$, ఒకే \tetoken{స్థిరసంకేతం}{!!{constant}}~$\Obj a$ను, అలాగే
$\lnot$, $\land$, $\lexists$ అనే తార్కిక సంకేతాలను మాత్రమే కలిగిన
ఒక సరళ మొదటిస్థాయి భాష కోసం ``\tetoken{సూత్రం}{!!{formula}}''ను నిర్వచిద్దాం. మన పూర్తి
నిర్వచనాలు మరింత సాధారణంగా ఉంటాయి: అనంత సంఖ్యలో \tetoken{విధేయ సంకేతాలు}{!!{predicate}s},
\tetoken{స్థిరసంకేతాలను}{!!{constant}s} అనుమతిస్తాం. నిజానికి, ``పదాలు'' రూపొందించడానికి
\tetoken{స్థిరసంకేతాలు}{!!{constant}s}, \tetoken{చరాలతో}{!!{variable}s} కలపగల \tetoken{ప్రమేయ సంకేతాలను}{!!{function}s} కూడా
పరిశీలిస్తాం. ప్రస్తుతానికి, $\Obj a$, చరరాశులు మాత్రమే మన పదాలు.
మనకు అనంత సంఖ్యలో \tetoken{చరాలు}{!!{variable}s} అవసరం. అధికారికంగా $\Obj v_0$,
$\Obj v_1$, \dots సంకేతాలను చరరాశులుగా వాడతాం.

\begin{defn}
\emph{\tetoken{సూత్రాలు}{!!{formula}s}} సమితి~$\Frm$ను కింది విధంగా నిర్వచిస్తాం:
\begin{enumerate}
\item\ollabel{fmls-atom} $\Atom{\Obj P}{\Obj a}$ మరియు $\Atom{\Obj
  P}{\Obj v_i}$లు \tetoken{సూత్రాలు}{!!{formula}s} ($i \in \Nat$).

\tagitem{prvNot}{\ollabel{fmls-not}$!A$ ఒక \tetoken{సూత్రం}{!!a{formula}} అయితే, $\lnot !A$
  ఒక \tetoken{సూత్రం}{!!{formula}}.}{}

\tagitem{prvAnd}{$!A$, $!B$లు \tetoken{సూత్రాలు}{!!{formula}s} అయితే, $(!A \land
  !B)$ ఒక \tetoken{సూత్రం}{!!a{formula}}.}{}

\tagitem{prvEx}{\ollabel{fmls-ex}$!A$ ఒక \tetoken{సూత్రం}{!!a{formula}}, $x$ ఒక \tetoken{చరం}{!!a{variable}}
  అయితే, $\lexists[x][!A]$ ఒక \tetoken{సూత్రం}{!!a{formula}}.}{}

\tagitem{limitClause}{\ollabel{fmls-limit}మరేదీ \tetoken{ఒక సూత్రం}{!!a{formula}} కాదు.}{}
\end{enumerate}
\end{defn}

ఏ $i \in \Nat$కైనా $\Atom{\Obj P}{\Obj a}$,
$\Atom{\Obj P}{\Obj v_i}$లు \tetoken{సూత్రాలు}{!!{formula}s} అని \olref{fmls-atom}
చెబుతుంది. వీటినే \emph{పరమాణు} \tetoken{సూత్రాలు}{!!{formula}s} అంటారు. మనకు
ప్రారంభించడానికి ఇవి ఆధారమిస్తాయి. ఇప్పటికే రూపొందించిన
సూత్రాల నుంచి కొత్త \tetoken{సూత్రాలను}{!!{formula}s} రూపొందించే పద్ధతులను ఇతర
ఉపవాక్యాలు ఇస్తాయి. ఉదాహరణకు, \olref{fmls-not} ప్రకారం $\lnot
\Atom{\Obj P}{\Obj v_2}$ ఒక \tetoken{సూత్రం}{!!a{formula}}; ఎందుకంటే
$\Atom{\Obj P}{\Obj v_2}$ ఇప్పటికే \olref{fmls-atom} వల్ల ఒక
\tetoken{సూత్రం}{!!a{formula}}. తరువాత \olref{fmls-ex} వల్ల
$\lexists[\Obj v_2][\lnot \Atom{\Obj P}{\Obj v_2}]$ మరో
\tetoken{ఒక సూత్రం}{!!{formula}}; ఇలాగే కొనసాగుతుంది. ఈ విధంగా మనం రూపొందించగల
సంకేతమాలలు \emph{మాత్రమే} \tetoken{సూత్రాలుగా}{!!{formula}s} లెక్కవుతాయని
\olref{fmls-limit} చెబుతుంది. ముఖ్యంగా,
$\lexists[\Obj v_0][\Atom{\Obj P}{\Obj a}]$,
$\lexists[\Obj v_0][\lexists[\Obj v_0][\Atom{\Obj P}{\Obj a}]]$
\tetoken{సూత్రాలుగా}{!!{formula}s} \emph{లెక్కవుతాయి}; బయట అనవసరమైన కుండలీకరణలు ఉన్నందున
$(\lnot \Atom{\Obj P}{\Obj a})$ మాత్రం లెక్కవదు.

\tetoken{సూత్రాలను}{!!{formula}s} ఇలా నిర్వచించడాన్ని \emph{ఆగమనాత్మక నిర్వచనం}
అంటారు. ఇది \emph{నిర్మాణాత్మక ఆగమనం} అనే ఆగమన నిరూపణ రూపాన్ని
ఉపయోగించి \tetoken{సూత్రాలు}{!!{formula}s} గురించి విషయాలను నిరూపించడానికి వీలు
కల్పిస్తుంది. వీటిని \olref[mth][ind][idf]{sec},
\olref[mth][ind][sti]{sec}లలో సాధారణంగా చర్చించాం; తరువాతి
నిరూపణల్లోకి వెళ్లే ముందు వాటిని పునశ్చరణ చేయాలి. ప్రాథమిక భావన ఇదీ:
ఏదైనా విషయం అన్ని \tetoken{సూత్రాలకు}{!!{formula}s} సత్యమని నిరూపించాలంటే, ముందుగా అది
పరమాణు \tetoken{సూత్రాలకు}{!!{formula}s} సత్యమని చూపాలి; తరువాత, అది ఏ
\tetoken{సూత్రం}{!!{formula}}~$!A$కు (మరియు~$!B$కు) సత్యమైతే, $\lnot !A$,
$(!A \land !B)$, $\lexists[x][!A]$లకు కూడా సత్యమని చూపాలి.
ఉదాహరణకు, అది $\lexists[\Obj v_2][\lnot \Atom{\Obj P}{\Obj
v_2}]$కు సత్యమని ఇది నిరూపిస్తుంది: మొదటి భాగం నుంచి పరమాణు
\tetoken{సూత్రం}{!!{formula}}~$\Atom{\Obj P}{\Obj v_2}$కు అది సత్యమని తెలుస్తుంది.
అప్పుడు రెండవ భాగం నుంచి $\lnot \Atom{\Obj P}{\Obj v_2}$కు సత్యమని,
మరొకసారి అదే భాగం నుంచి $\lexists[\Obj v_2][\lnot \Atom{\Obj
P}{\Obj v_2}]$కే సత్యమని లభిస్తుంది. అన్ని \tetoken{సూత్రాలు}{!!{formula}s} పరమాణు
\tetoken{సూత్రాలు}{!!{formula}s} నుంచి ఆగమనాత్మకంగా ఉత్పన్నమవుతాయి కాబట్టి, వాటిలో
దేనికైనా ఇది పనిచేస్తుంది.

\end{document}
